\documentclass[sigplan, nonacm = true, natbib = false]{acmart}

\usepackage{listings}

\lstset{
  basicstyle=\ttfamily\footnotesize,
  columns=fullflexible,
  keepspaces=true
}

\lstdefinestyle{metl}{
  language=[Objective]Caml,
  morekeywords={eff, handle, end, return, mask, forall, Effect, alias},
  morecomment=[l]{--},
  literate=
    {->}{{$\to$}}2
    {=>}{{$\Rightarrow$}}2
    {()}{{()}}2
}

\lstdefinestyle{frank}{
  language=Haskell,
  morekeywords={interface},
  literate=
    {->}{{$\to$}}2
}

\lstdefinestyle{ocaml}{
  language=[Objective]Caml,
  morekeywords={},
  literate=
    {->}{{$\to$}}2
    {=>}{{$\Rightarrow$}}2
    {()}{{()}}2
}

\lstdefinestyle{koka}{
  morekeywords={fun,rec,handler,forall, with, return, ctl, effect, div, type},
  sensitive=true,
  morecomment=[l]{--},
  morestring=[b]",
  literate=
    {->}{{$\to$}}2
    {()}{{()}}2
}

\usepackage[backend=bibtex, style=numeric]{biblatex}
\addbibresource{abstract.bib}
\renewcommand*{\bibfont}{\footnotesize}

%% Keywords

\newcommand{\keyw}[1]{{\textbf{{\textsf{#1}}}}}
\newcommand{\Do}{\keyw{do}\;}
%%

\title{Higher-order fork, modaly (Extended Abstract)}

\settopmatter{authorsperrow=3}

\author{Aghilas Y. Boussaa}
\email{aghilas.boussaa@ens.fr}
\affiliation{%
  \institution{École normale supérieure -- PSL}
  \country{France}
}

\author{Wenhao Tang}
\orcid{0009-0000-6589-3821}
\email{wenhao.tang@ed.ac.uk}
\affiliation{%
  \institution{The University of Edinburgh}
  \country{UK}
}

\author{Sam Lindley}
\orcid{0000-0002-1360-4714}
\email{sam.lindley@ed.ac.uk}
\affiliation{%
  \institution{The University of Edinburgh}
  \country{UK}}

\keywords{effect handlers, effect types, modal types}

\begin{document}
\sloppy
\begin{abstract}
  Effect handlers are a programming construct allowing users to define and
  combine computational effects. An effect of interest is cooperative
  concurrency. Various static type-and-effect systems can guarantee
  safety. Implementing it with a higher-order fork can show shortcomings of such
  a system.

  We review implementation of cooperative concurrency in existing languages, and
  in a new one based on modal effect types.
\end{abstract}
\maketitle
\section{Introduction}
Effect handlers can express many control-flow constructs including cooperative
concurrency~\cite{structuredasync2017}. One can model it with an effect with two
operations \textsf{fork} and \textsf{yield} that, respectively create a new task
and pause the current one. The languages we will consider here
(\textsf{Koka}~\cite{structuredasync2017, koka2014},
\textsf{Frank}~\cite{frank2020}, \textsf{Effekt}~\cite[Example~10]{effekt2022},
\textsf{\textsc{Metl}}~\cite[Section~2.6]{tang2025modal}) can implement
multitasking with a UNIX-style~\cite{hillerstrom2021thesis} \textsf{fork} which
has type: $\textsf{fork} : \textsf{unit} \Rightarrow \textsf{bool}$. The
continuation receives a boolean indicating whether it is the parent or child,
like in the following piece of \textsf{\textsc{Metl}} code.
 gs\begin{lstlisting}[style=metl]
  if do fork () then
    print "parent"
  else
    print "child"
\end{lstlisting}
However, this approach relies on \emph{multi-shot} continuations, while some
existing languages that gain support for effect handlers only support
\emph{single-shot} continuations~\cite{ocamleffect2021, wasmeffect2023}.

One can still implement cooperative concurrency with effect handlers in
\textsf{OCaml} which only supports single-shot
continuation~\cite{dolan2015effective, ocamleffect2021}. The \textsf{fork}
operation has then the following signature.
\begin{lstlisting}[style=ocaml]
  type _ Effect.t += Fork : (unit -> unit) -> unit t
\end{lstlisting}
This operation takes as input a thunk of the computation which the new task
should run. The scheduler maintains a queue of thunks that it runs. As
\textsf{OCaml} has no static effect system, the implementation and the type
signatures are straightforward.
\section{Implementation in existing languages}
We will now give an implementation of a cooperative concurrency with a
\textsf{fork} that takes a thunk as input in several languages, focusing on the
type signatures and some details of implementation. To demonstrate
compositionality of the type-and-effect systems, we will suppose the existence
of a parametric \textsf{queue} effect with two operations: \textsf{enqueue} and
\textsf{dequeue}, and use it to store thunks.

All the implementations are in appendix and available \href{https://github.com/4y8/hope26}{online}.
\subsection{\textsf{Koka}}
\textsf{Koka}~\cite{koka2014} has a traditional type-and-effect system where
each arrow type holds a scoped row~\cite{scopedrows2005} of effects the function
can perform. We rely on effect polymorphism to store and yield effectful
functions while pleasing the type system, modeling concurrency with this effect.
\begin{lstlisting}[style=koka]
  div effect co<e>
    ctl yield () : unit
    ctl fork (p : () -> <co<e>,div|e> unit) : unit
\end{lstlisting}
\textsf{Koka}'s positivity checker requires the \textsf{div} keyword before the
declaration, and we propagate the \textsf{div} (which signals potential
non-termination) effect throught all of our signatures. The type of processes is
also effect-polymorphic:
\begin{lstlisting}[style=koka]
  div type proc<e>
    Proc (p : () -> <queue<proc<e>>,div|e> unit).
\end{lstlisting}
With help from the alias \lstinline[style = koka]{qproc e = queue<proc<e>>}, we
can implement a scheduler with type:
\begin{lstlisting}[style=koka]
  (() -> <co<e>, qproc<e>, div|e> unit) ->
    <qproc<e>,div|e> unit,
\end{lstlisting}
which becomes 
\begin{lstlisting}[style=koka]
  (() -> <co<e>, qproc<e>|e> unit) -> <qproc<e>|e> unit,
\end{lstlisting}
if we forget the \textsf{div} effect.
\subsection{\textsf{Frank}}
Noticing that the trailing effect variable is often the same through a
signature, the \textsf{Frank}~\cite{frank2020} language gives syntactic sugar
allowing users to omit them.

The paper cited above has a concurrency example in Section~6 which we transcribe
here. The corresponding effect has signature:
\begin{lstlisting}[style=Frank]
  interface Co = fork : {[Co]Unit} -> Unit
               | yield : Unit
\end{lstlisting}
where \lstinline[style=Frank]!{[Co]Unit}! is the type of computations with
effect \lstinline[style=Frank]{Co}. By defining the type
\lstinline[style=Frank]{Proc} to be \lstinline[style=Frank]!proc {[Queue Proc]Unit}!,
we can implement a scheduler with the following variable-free type.
\begin{lstlisting}[style=Frank]
  {<Co>Unit -> [Queue Proc]Unit}
\end{lstlisting}
The type checker desugars this type to:
\begin{lstlisting}[style = Frank]
  {<Co e>Unit -> [Queue (Proc e), e]Unit}.
\end{lstlisting}
\subsection{\textsf{Effekt}}
We have not managed yet to implement our example in
\textsf{Effekt}~\cite{effekt2020, effekt2022}.
\subsection{\textsf{\textsc{Metl}}}
Drawing inspiration from Multimodal Type Theory~\cite{mtt2020}, modal effect
types~\cite{tang2025modal, tang2026modal} use \emph{modalities} to give a robust
version of \textsf{Frank}'s idea of \emph{ambient ability}, called the
\emph{ambient effect context}, eliminating effect variables. To guarantee effect
safety, operations can only yield values whose meaning is independant from the
ambient effect context, which is not the case of arrow types, as they may use
any effect available. Yielding a function is only possible by first boxing it in
an \emph{absolute} modality that fixes the set of effects it can
perform. However, to use the \textsf{co} effect with others, we cannot fix a set
of effects and need to bring back effect
variables~\cite[Section~2.10]{tang2025modal}.
\begin{lstlisting}[style=metl]
  eff co (e : Effect)
    = fork : [co e, e](unit -> unit) => unit
    | yield : unit => unit
\end{lstlisting}
We define the polymorphic type $\textsf{proc}$ of processes like in
\textsf{Koka}.
\begin{lstlisting}[style=metl]
  type proc (e : Effect) =
    Proc of [queue (proc e), e](unit -> unit).
\end{lstlisting}
Then, the \textsf{scheduler} has type:
\begin{lstlisting}[style=metl]
  forall (e : Effect) . [queue (proc e), e]
    ([co e, queue (proc e), e](unit -> unit) -> unit),
\end{lstlisting}

This type signature reminds the \textsf{Koka} one, erasing all the benefits of
modal effect types, as each arrow comes with an absolute modality. Moreover, the
implementation relies on a \emph{modality-parametrised
handler}~\cite[Section~3.5]{tang2026modal} to give a type with an absolute
modality to the continuation bringing extra burden compared to \textsf{Koka}.
\section{Higher-order effects in \textsf{\textsc{Met}}}
The unsatisfying solution in \textsf{\textsc{Metl}} motivates the search for
other ways of extending the core calculus to support our example. We have
implemented a promising approach: relax kind restrictions on operations’
operands, but restrict subeffecting in presence of non-absolute types. These
conditions should guarantee that a non-absolute value arrives to an effect
context more permessive than the one it comes from.

With these new rules, we can define the \textsf{co} effect as follows.
\begin{lstlisting}[style=metl]
  eff! co = fork : (unit -> unit) => unit
          | yield : unit => unit
\end{lstlisting}
The bang indicates that \textsf{co} is a higher-order operation. Unlike in
previous examples, the signature of \textsf{fork} does not mention the
\textsf{co} effect. In fact, the thunk will have the same ambient effect context
as the computation that yielded it, which contains effect \textsf{co}. For the
same reason, we do not need a distinct type of processes and can use
functions. The scheduler has then the following type, and it does not rely on
modality-parametrised handler.
\begin{lstlisting}[style=metl]
  [queue (unit -> unit)](<co>(unit -> unit) -> unit)
\end{lstlisting}
Another difference with the other implementations is the absence of masking in
the scheduler's implementation.

For now, we lack a proof of safety and an expressivity result relating this
system with effect polymorphism.

\printbibliography
\appendix
\section{\textsf{Koka}}
\lstinputlisting[style=koka]{impl/koka.kk}
\section{\textsf{Frank}}
\lstinputlisting[style=frank]{impl/frank.fk}
\section{\textsf{\textsc{Metl}} with effect variables}
\lstinputlisting[style=metl]{impl/metl-eff-var.met}
\section{\textsf{\textsc{Metl}} with higher-order effects}
\lstinputlisting[style=metl]{impl/metl-ho-eff.met}
\end{document}
